\chapter{Минимальный пакет новых данных dynamic parent charged mediator}
\label{ch:v8-charged-mediator-minimal-new-data}

\section{Один резонансный квант заменяет два масштаба}

Предыдущий ledger перечислял отдельно mediator-gap $7\gamma$ и clock-energy
$E_C$. Энергосохраняющий collision требует их совпадения, поэтому введём
один положительный квант
\begin{equation}
 E_*=E_C=7\gamma>0.
 \label{eq:v8-mediator-min-data-common-energy}
\end{equation}
Тогда прежний параметр автоматически равен
\begin{equation}
 \gamma=\frac{E_*}{7}.
 \label{eq:v8-mediator-min-data-gamma}
\end{equation}
Это типизированное сокращение данных, а не численный selector величины
$E_*$.

\section{Coupling и tick выражаются через два непрерывных входа}

Пусть $\chi>0$ есть безразмерная относительная сила interaction. Его
размерная амплитуда равна
\begin{equation}
 g=\chi E_*.
 \label{eq:v8-mediator-min-data-coupling}
\end{equation}
Clock tick и weak-limit rate после этого определяются формулами
\begin{equation}
 \tau_C=\frac{\hbar}{E_*},
 \qquad
 \Gamma=\chi^2\frac{E_*}{\hbar}.
 \label{eq:v8-mediator-min-data-tick-rate}
\end{equation}
Они удовлетворяют двум точным связям
\begin{equation}
 E_*\tau_C=\hbar,
 \qquad \Gamma\tau_C=\chi^2,
 \label{eq:v8-mediator-min-data-relations}
\end{equation}
а collision-формула записывается без нового параметра:
\begin{equation}
 \Gamma=\frac{g^2\tau_C}{\hbar^2}.
 \label{eq:v8-mediator-min-data-collision-rate}
\end{equation}

\section{Два непрерывных параметра необходимы}

Рассмотрим карту
\begin{equation}
 F(E_*,\chi)=
 \left(\frac{E_*}{7},E_*,\chi E_*,
 \frac{\hbar}{E_*},\frac{\chi^2E_*}{\hbar}\right).
 \label{eq:v8-mediator-min-data-map}
\end{equation}
Её якобиан имеет ранг два, поскольку minor по координатам $(E_C,g)$ равен
\begin{equation}
 \det\frac{\partial(E_C,g)}{\partial(E_*,\chi)}=E_*>0.
 \label{eq:v8-mediator-min-data-jacobian-minor}
\end{equation}
Следовательно, один непрерывный параметр не может независимо задавать общий
масштаб и относительную силу.

Два точных scale-witness при $\hbar=1$ дают
\begin{equation}
 (E_*,\chi)=(1,1),(2,1)
 \quad\Longrightarrow\quad \Gamma_2/\Gamma_1=2,
 \label{eq:v8-mediator-min-data-scale-witness}
\end{equation}
а coupling-witnesses
\begin{equation}
 (E_*,\chi)=(1,1),(1,2)
 \quad\Longrightarrow\quad \Gamma_2/\Gamma_1=4.
 \label{eq:v8-mediator-min-data-coupling-witness}
\end{equation}
Обе свободы сохраняют тип carrier и законы сохранения.

\section{Полный минимальный пакет}

Кроме двух непрерывных координат остаются два независимых структурных
выбора: само endpoint/carrier extension и класс транспортного закона.
Итак,
\begin{equation}
 \mathcal D_{\min}=
 \{\mathfrak e_{\rm endpoint},E_*,\chi,\mathfrak t_{\rm transport}\},
 \qquad |\mathcal D_{\min}|=4.
 \label{eq:v8-mediator-min-data-package}
\end{equation}
Вакуум не является пятым независимым входом: после задания положительного
локального parent он есть его единственное нулевое состояние,
\begin{equation}
 \ker h_{45}=\mathbb C|0\rangle.
 \label{eq:v8-mediator-min-data-vacuum-derived}
\end{equation}
Поэтому предыдущие пять apparent slots сокращаются на один:
\begin{equation}
 N_{\rm apparent}=5,
 \qquad N_{\rm independent}=4.
 \label{eq:v8-mediator-min-data-reduction}
\end{equation}
Точная dependency-классификация закрыта,
\begin{equation}
 N_{\rm dependency}=7/7,
 \qquad N_{\rm minimal\ classification}=4/4,
 \label{eq:v8-mediator-min-data-ledgers}
\end{equation}
но текущий parent не выбирает ни одного из четырёх слотов:
\begin{equation}
 N_{\rm inherited\ selection}=0/4.
 \label{eq:v8-mediator-min-data-origin-ledger}
\end{equation}

\begin{theorem}[Минимальность четырёхслотового пакета]
После резонансной идентификации $E_*=E_C=7\gamma$ все размерные динамические
величины параметризуются двумя и только двумя непрерывными координатами
$(E_*,\chi)$. Вместе с endpoint-extension и transport primitive это даёт
ровно четыре независимых новых parent-данных. Ни одно из них текущим
действием не выбрано.
\end{theorem}

\begin{nogocriterion}[Зависимость не является выводом значения]
Равенства $\gamma=E_*/7$, $g=\chi E_*$ и
$\tau_C=\hbar/E_*$ устраняют избыточные координаты, но не фиксируют
$E_*$ или $\chi$. Нельзя объявлять сокращение числа параметров физическим
предсказанием их значений.
\end{nogocriterion}

\begin{protocol}\begin{enumerate}
\item общий резонансный квант: \textbf{$E_*=E_C=7\gamma$};
\item независимых непрерывных параметров: \textbf{$2$};
\item их координаты: \textbf{$(E_*,\chi)$};
\item ранг dependency-Jacobian: \textbf{$2$};
\item размерный coupling: \textbf{$g=\chi E_*$};
\item tick: \textbf{$\tau_C=\hbar/E_*$};
\item rate: \textbf{$\Gamma=\chi^2E_*/\hbar$};
\item независимых structural slots: \textbf{$2$};
\item полный минимальный пакет: \textbf{$4$ слота};
\item dependency reduction: \textbf{$7/7$};
\item inherited selection: \textbf{$0/4$};
\item следующий гейт: \textbf{parent-origin общего энергетического масштаба}.
\end{enumerate}\end{protocol}

Машинный сертификат сохранён в
\path{s2t_v8_baryon_c0_charged_mediator_dynamic_parent_minimal_new_data_gate_results.json}.

\begin{thebibliography}{9}
\bibitem{v8-mediator-min-data-attal} S.~Attal and Y.~Pautrat,
\emph{From Repeated to Continuous Quantum Interactions},
Ann. Henri Poincar\'e 7 (2006) 59; arXiv:math-ph/0311002.
\end{thebibliography}